\begin{lemma}[Proved: primary parity for the genuine mod-$p^2$ Tate Bockstein]
\label{C:modp2_Tate_Bockstein}
Let $A$ be a characteristic-zero DVR with maximal ideal $(p)$ and
residue field $F$. For a finite free $A/p^2$-representation $M$ of
$C_p$, use the two-periodic complex
\[
 d^0=T-I,\qquad d^1=I+T+\cdots+T^{p-1}.
\]
On its homology modulo $p$ there is a natural odd Bockstein $b$:
if $\bar v$ is a cycle, lift it and write $dv=pw$, then put
$b[\bar v]=[\bar w]$. For every monic irreducible
$q\in F[z]$ other than $z$, the function
\[
 \chi_q(M)=\operatorname{length}_{F[z]}
     \bigl(H^0(M/pM),z=b^2\bigr)_q\pmod2
\]
is additive on admissible exact sequences. It is zero on every
representation obtained by reduction from a finite free
$A$-representation of $C_p$.
\end{lemma}
\begin{proof}
Both products of the two differentials are zero, since $T^p=I$.
For a free $A/p^2$-module the kernel of multiplication by $p$ is
its image, so division by $p$ determines $\bar w$ uniquely.
The equality $0=d^2v=p\,dw$ implies that $\bar w$ is a cycle.
Replacing $v$ by $v+pu$ changes $\bar w$ by the boundary $d\bar u$;
if $\bar v$ is a boundary, choose the lift $v=dz$, which gives zero.
This proves that $b$ is well defined; the same calculation with a
chain map proves naturality.

For an admissible sequence, the underlying free modules split.
Write the differential on its middle term as
$\left(\begin{smallmatrix}d_A^i&X^i\\0&d_B^i\end{smallmatrix}\right)$.
Its square-zero identity says
$d_A^{i+1}X^i+X^{i+1}d_B^i=0$.
If $d_B^iv=pw$, division by $p$ gives, on homology,
$b_A\partial+\partial b_B=0$.
The other two maps of the Tate six-cycle commute with $b$.
Consequently the six-cycle respects the primary decomposition for
$z=b^2$. On every part with $q\ne z$, $b$ is invertible between
the even and odd groups. It identifies the six-cycle with an exact
cyclic sequence on three even groups. More explicitly, if the odd
groups are identified using $b^0$, the other connecting map differs
from the folded one by the unit $-z$; this preserves the required
kernels and images, since every map commutes with $z$. Its lengths are
$r+t,r+s,s+t$ for the three image lengths $r,s,t$.
Their sum is even, proving additivity.

Finally, if the representation lifts to $A$, its Tate differential
lifts to a square-zero differential over $A$. Given a lift $v$ of
a mod-$p$ cycle, write $dv=pw$ over $A$. Since multiplication by
$p$ is injective on the free module, $d^2v=0$ implies $dw=0$.
Thus $b^2=0$ on mod-$p$ homology. All the indicated nonzero primary
parts vanish, as claimed.
\end{proof}

\begin{proposition}[Proved: height-zero odd-primary torsion is an obstruction]
\label{C:odd_torsion_heightzero}
Let $p$ be odd, let $A$ be as in the preceding lemma, and let $J$
be any $A$-module. If
\[
 J[p]\not\subseteq pJ,
\]
then $\vartheta_{A\ltimes J}$ is not surjective in degree $p$.
In particular this holds for $\mathbb Z_{(p)}[\epsilon]/
(p\epsilon,\epsilon^2)$, for every odd prime $p$.
The obstruction can be represented by the difference of two
rank-$(p-1)!$ induced representations and vanishes after inverting
$p$ and at every residue-field point.
\end{proposition}
\begin{proof}
Choose $x\in J[p]$ with nonzero class in $J/pJ$, and an $F$-linear
functional $\ell:J/pJ\to F$ with $\ell(\bar x)=1$.
The ring $A\ltimes J$ is local, so $K_0(A\ltimes J)=\mathbb Z$;
its degree-$p$ polynomial image is generated by the Young
permutation classes.
Identify $F$ with $pA/p^2A$ by multiplication by $p$.
Then
\[
 A\ltimes J\longrightarrow A/p^2,
       \qquad (a,j)\longmapsto a+p\ell(\bar j)
\]
is a well-defined ring homomorphism: the second summand is
$A$-linear and its products are zero. In particular it sends $x$
to $p$.

Since $J^2=0$ and $px=0$, the free line $L_x$ on which a chosen
generator of $C_p$ acts by $1+x$ is a representation.
After the displayed scalar extension its generator is $1+p$.
Its two Tate differentials are $p$ and $p$ modulo $p^2$, since
\[
 \sum_{i=0}^{p-1}(1+p)^i
       =p+p^2\frac{p-1}{2}=p\pmod{p^2}.
\]
Its mod-$p$ differential is zero; both Bocksteins are one, and
$b^2=1$. Under the twist $g\mapsto g^r$, $1\le r<p$, the two
Bocksteins are $r$ and $1$. Exactly the twist $r=1$ therefore has
nonzero $(z-1)$-primary parity.

Put $G=\mathfrak S_p$, $P=C_p$, and
$\xi=[\operatorname{Ind}_P^G L_x]-[\operatorname{Ind}_P^G\mathbf1]$.
Mackey restriction to $P$ contributes the $p-1$ normalizer twists
once each; every remaining double coset has trivial intersection
and gives an induced free $A/p^2[P]$-module, whose mod-$p$ Tate
complex is acyclic. Lemma~\ref{C:modp2_Tate_Bockstein} with
$q=z-1$ therefore evaluates to one on $\xi$.
Every Young permutation representation lifts to $A$, so its value
is zero. Hence $\xi$ is outside the degree-$p$ polynomial image.
The ranks are $[G:P]=(p-1)!$. Finally $x$ becomes zero after
inverting $p$, and it maps to zero in every residue field because
it is square-zero. The two representations then agree.
\end{proof}

This statement makes no assertion when $J[p]\subseteq pJ$,
notably for $p$-divisible torsion modules. The earlier flat-module
criterion and this proposition do not by themselves classify all
square-zero $A$-modules at odd primes.

The finite certificate \texttt{code/C\_finish\_modp2\_Bockstein.py} tests all $5051$ stated extension candidates and the $927$ admissible ones at $p=3,5,7$, computing actual mod-$p$ Tate homology, divided differentials, connecting maps, and primary parity. Output is in \texttt{logs/C\_finish\_modp2\_Bockstein.json}.
